\section{Introduction}
\label{sec:introduction}

CMS Computing Operations
(CompOps)~\cite{cms_computing_tdr,cms_compops} turns raw CMS
detector output into usable scientific data through three
activities (Section~\ref{sec:workload}): Tier-0 repacks the
detector stream, archives two tape copies, and runs express and
prompt reconstruction; Production \& Reprocessing drives Monte
Carlo and reprocessing campaigns through the CMS Workflow
Management system; and Data
Management~\cite{cms_dm2025} executes transfers and placement
across more than sixty Tier-1 and Tier-2 WLCG sites with
Rucio~\cite{rucio2014,cms_rucio2020}.
Day-to-day, CompOps operators run workflows, triage tickets,
watch live operational state (the MonIT index, surfaced through
OpenSearch and Grafana), conduct root-cause analysis, and onboard
new team members; their support is fragmented across the CMS
TWiki, the CompOps Jira projects, MonIT, and Mattermost channels,
and resolving a stalled transfer requires the runbook entry, the
current MonIT counters, and the ticket history of past stalls on
the same link.

Archi is an operator chat agent deployed for CompOps since 16
February 2026. Archi combines retrieval over the CMS TWiki and
CompOps Jira with live MonIT calls and surfaces the tool trace
for operator audit. The production pipeline uses GPT-5; to
evaluate local-weight alternatives we sample a 260-question
subset from the recorded conversations after deduplication and
filtering (Section~\ref{sec:question_patterns}; 144 answerable
from documentation alone, 116 requiring a live tool call) and
re-score the GPT-5.3 production answers on the same rubric as
the GPT-5.3 reference.

Each pipeline layer earns a measurable lift on the rubric. On
the largest open-weight model the bare LLM scores 2.86, the
RAG control 3.97, the agent without live tools 4.22, and the
full agent 4.22 in aggregate with $+0.19$ on the 20\% of
operator traffic that needs operational state. The full agent
at $4.22 \pm 0.06$ matches the GPT-5.3 production reference at
$4.08 \pm 0.08$ under the same judge.
Section~\ref{sec:agent} describes the agent and its tool
surface, Section~\ref{sec:evaluation} the four-layer rubric
evaluation and a four-judge replication, and
Section~\ref{sec:discussion} the trade-offs the agent loop
carries.
